% !TeX program = xelatex
\documentclass[12pt,letterpaper]{article}
\usepackage{iftex}
\ifPDFTeX
  \usepackage[T1]{fontenc}
  \usepackage{newtxtext}
\else
  \usepackage{fontspec}
  \setmainfont{Times New Roman}
\fi
\usepackage[hidelinks]{hyperref}
\usepackage[letterpaper,top=0.75in,bottom=0.75in,left=0.75in,right=0.75in,headheight=14.5pt,headsep=3pt,footskip=18pt]{geometry}
\usepackage{fancyhdr}
\usepackage{needspace}
\pagestyle{fancy}
\fancyhf{}
\fancyhead[R]{Hamza Abdelhedi}
\fancyfoot[C]{\thepage}
\renewcommand{\headrulewidth}{0pt}
\renewcommand{\footrulewidth}{0pt}
\AtBeginDocument{\fontsize{12pt}{12pt}\selectfont}
\setlength{\parindent}{0pt}
\setlength{\parskip}{0pt}
\setlength{\emergencystretch}{3em}
\raggedbottom
\newcommand{\proposalheading}[1]{\begin{center}\textbf{#1}\end{center}\vspace{-4pt}}
\newcommand{\nsercsection}[1]{\Needspace{2\baselineskip}\vspace{3pt}\noindent\textbf{#1}\par\vspace{0pt}}
\newcommand{\citep}[1]{[#1]}
\begin{document}
\proposalheading{Brain-wide neural dynamics of human decisions under changing evidence}

\nsercsection{Background and rationale}
Perceptual decisions are often studied while evidence remains stationary within a trial. These paradigms have established influential accumulation-to-bound accounts, but distinct termination mechanisms---including fixed or collapsing bounds and urgency-gating---can produce overlapping behavioural predictions under some conditions \citep{1--3}. Dynamic decisions provide stronger leverage because the evidence itself changes while deliberation is ongoing. The brain must balance sensitivity to new information against the cost of waiting, determine how much earlier evidence should still constrain the current state, and eventually commit. The Tokens task was designed around this problem: sequential evidence changes within each trial and participants may respond before all information is revealed \citep{1}. In monkeys, premotor and motor populations track changing evidence together with urgency and show a neural signature of commitment before movement \citep{4}; broader analyses indicate that decision formation emerges from coordinated dynamics across multiple cortical and subcortical areas \citep{5,6}. Human EEG studies likewise separate evidence-, decision-variable-, and urgency-related signals \citep{7}, and show that temporal evidence weighting adapts to environmental statistics \citep{8} and can be influenced by confidence-dependent control \citep{9}. Recent invasive work further suggests that commitment may involve a transition from input-driven to more autonomous population dynamics \citep{10}.

The remaining gap is therefore not whether dynamic decisions have neural correlates, but how distributed human neural states implement the evolving computation. In particular, we do not know how current evidence, recent evidence history, urgency, and large-scale coordination jointly organize brain-wide activity; how that state is redirected when new evidence contradicts an emerging choice; or whether commitment reflects continuous evolution or a distinct transition that can still be revised. These questions require within-trial evidence changes together with millisecond-scale, whole-brain measurements.

\nsercsection{Preliminary results}
The primary dataset provides exactly this combination: 275-sensor MEG with individual MRI from 32 healthy adults performing the Tokens task (16,324 trials). Thiery et al. previously showed that these data contain low-dimensional oscillatory trajectories in which alpha- and beta-band components track evolving sensory information across visual and sensorimotor systems, with alpha activity related to choice commitment and beta activity to motor execution \citep{11}. My ongoing analyses reproduce these core trajectory results, and I am leading a first-author replication and extension that adds source-resolved analyses of large-scale functional coordination. Behaviourally, mean accuracy is 83.7\%, and decisions are faster in Fast than Slow blocks (approximately 1186 vs. 1314 ms), confirming robust adaptation to temporal context. Psychophysical kernels show systematic temporal weighting of the token sequence. I have also compared four sequential-sampling mechanisms using held-out likelihood; among the models tested, a multiplicative urgency account is currently preferred in 28/32 participants, with model identity recovered in 47/48 simulations. I treat this as a quantitative constraint, not proof of a biological mechanism: competing models will generate trial-wise predictions for neural adjudication. Critically, approximately 10.5\% of trials are initially misleading, so later evidence opposes the emerging choice. Behaviour on these trials shows that later evidence can correct an early tendency, providing a natural within-trial perturbation for testing neural updating and commitment.

\nsercsection{Central hypothesis and objectives}
I hypothesize that dynamic decisions emerge from evolving brain-wide neural states jointly shaped by current sensory evidence, recent evidence history, and time-dependent commitment pressure. Controlled changes and reversals in evidence will reveal how strongly prior information constrains these states, while late contradictory evidence will determine whether commitment behaves as an effectively absorbing endpoint or remains revisable.

\textbf{Aim 1 --- Determine how evidence, urgency, and distributed coordination organize brain-wide decision dynamics.} Building on the replicated oscillatory trajectories, I will quantify how momentary evidence, temporally weighted evidence history, elapsed time/urgency, and proximity to commitment explain the evolving source-resolved neural state. Rather than asking only where these variables are represented, I will test whether they occupy distinguishable dimensions of the distributed state space and how their relative contributions change as a decision forms. Source-level connectivity will determine whether commitment is accompanied by reproducible reconfiguration of large-scale coordination, extending the trajectory analysis from distributed activity to distributed interaction. I predict that evidence-related and urgency-related information will explain partially separable components of the neural state and that the transition toward commitment will be accompanied by systematic reorganization of inter-regional coordination. A small set of complementary complexity measures will be used only where informative to characterize broader changes in global brain state.

\textbf{Aim 2 --- Determine how brain-wide neural states update when evidence changes or reverses.} Misleading trials provide a direct perturbation of the evolving decision. I will estimate neural evidence-history kernels: at each moment, the influence of evidence received at successive earlier moments on the current neural state. I will combine them with trajectory geometry to quantify the direction and speed with which the state is redirected after evidence changes sign, and compare Fast and Slow blocks to test whether temporal weighting adapts to the cost of waiting. The key alternatives are mechanistically distinct: strong persistence predicts prolonged influence of earlier evidence and slower redirection after reversal, whereas stronger discounting predicts shorter history dependence and faster updating. I predict that the influence of past evidence will decay with lag and that its effective timescale will shorten when commitment pressure is greater, consistent with the broader principle that integration timescales can adapt to dynamic environments \citep{8,12}.

\textbf{Aim 3 --- Determine whether commitment corresponds to a transition in brain-wide dynamics and whether that state remains reversible.} Recent invasive work proposes that commitment coincides with a change from input-driven to internally dominated dynamics \citep{10}, while human work shows that decision termination rules can flexibly adjust with task demands \citep{13}. I will test whether an analogous transition is detectable at the whole-brain level in humans using convergent state-space and change-point analyses. Candidate transition times will be validated behaviourally using the psychophysical kernel: evidence arriving after a genuine commitment transition should exert less influence on choice than evidence arriving before it. Dynamic evidence then enables a stronger test than stationary trials. On late misleading trials, if commitment is effectively absorbing, contradictory post-transition evidence should have little behavioural influence and neural activity should remain in the committed regime. If commitment remains revisable, contradictory evidence should redirect the neural state toward deliberation or the alternative choice and predict behavioural changes of mind. Either outcome constrains the mechanism without assuming a particular state-transition model.

\nsercsection{Approach, generalization, and individual differences}
MRI-informed source reconstruction will provide cortical time series for oscillatory state-space, connectivity, and encoding analyses. Trial-wise latent quantities from competing decision models will be related to neural dynamics using cross-validation, matched surrogate controls, and across-participant stability tests. The claims concern distributed brain-wide dynamics, acknowledging the spatial limits and leakage structure of MEG source estimation; connectivity will therefore be interpreted as large-scale coordination rather than direct anatomical coupling.

An independent cohort of 74 healthy adults performing a dynamic face-morph task with 256-channel EEG and individual MRI provides a complementary test of generalization across task and modality. Published analyses in this cohort identified separable evidence-, decision-variable-, and urgency-related EEG signals, substantial between-participant variation in urgency, and 42 early versus 32 late responders \citep{7}. My preliminary EEG reanalysis reproduces this behavioural separation and is extending it to low-dimensional neural dynamics. I will test whether the evidence-, urgency-, and commitment-related organization identified in MEG generalizes to EEG and whether variation in these dynamics explains continuous differences in commitment strategy. Because urgency has been proposed to vary as an individual trait \citep{7,14}, independently measured behavioural or personality traits, where available, will provide an exploratory test of whether stable characteristics relate to neural-computational decision profiles; group labels will remain descriptive summaries of underlying continuous variation.

\nsercsection{Significance and feasibility}
This project will move human dynamic-decision neuroscience from describing decision-related signals to explaining how distributed neural states update, coordinate, and commit while evidence itself changes. It will provide mechanistic tests of evidence-history dependence, network reconfiguration, and commitment reversibility at a brain-wide scale, while determining which principles generalize across two dynamic-decision paradigms and individuals. Feasibility is unusually high: both datasets are already collected; the core MEG trajectory results are replicated in my ongoing analyses; behavioural model comparison, recovery, psychophysical kernels, and reversal-relevant trial definitions are established; and source-level analyses are operational. The resulting first-author MEG replication and connectivity extension will be submitted as part of this doctoral work, with analysis code released openly.
\end{document}
